\documentclass[11pt]{article}

% -------------------------------------------------
% Packages
% -------------------------------------------------
\usepackage[margin=1in]{geometry}

\usepackage{amsmath,amssymb,amsfonts,amsthm,mathtools,bm}
\usepackage{booktabs}
\usepackage{array}
\usepackage{multirow}
\usepackage{longtable}
\usepackage{enumitem}
\usepackage{xcolor}
\usepackage{graphicx}
\usepackage{float}
\usepackage{hyperref}
\usepackage{titlesec}
\usepackage{setspace}
\usepackage{microtype}
\usepackage{parskip}
\usepackage{tabularx}


% -------------------------------------------------
% Formatting
% -------------------------------------------------
\setstretch{1.1}

\titleformat{\section}
{\large\bfseries}
{\thesection.}{0.5em}{}

\titleformat{\subsection}
{\normalsize\bfseries}
{\thesubsection.}{0.5em}{}

\titleformat{\subsubsection}
{\normalsize\itshape}
{\thesubsubsection.}{0.5em}{}

\setlist[itemize]{topsep=2pt,itemsep=2pt}
\setlist[enumerate]{topsep=2pt,itemsep=2pt}

\hypersetup{
    colorlinks=true,
    linkcolor=blue,
    citecolor=blue,
    urlcolor=blue
}

% -------------------------------------------------
% Commands
% -------------------------------------------------
\newcommand{\E}{\mathbb{E}}
\newcommand{\Pbb}{\mathbb{P}}
\newcommand{\Var}{\mathrm{Var}}

% -------------------------------------------------
% Title
% -------------------------------------------------
\title{\textbf{Reference Notes}}
\author{Jonathan Ma}
\date{\today}

\begin{document}

\maketitle
\tableofcontents
\newpage

% =================================================================================
\part{Readings \& Literature Review}

\section{Age-Related Macular Degeneration: A Review}

\subsection*{Overview}

Age-related macular degeneration (AMD) is a progressive degenerative disease of the macula, the central retinal region responsible for high-acuity vision. AMD is among the leading causes of irreversible vision loss in adults over age 55 and is projected to affect nearly 288 million individuals worldwide by 2040.

The disease primarily affects photoreceptors, retinal pigment epithelium (RPE), bruch membrane, and choriocapillaris.
AMD progression involves degeneration of these structures, ultimately leading to central vision loss. Early and intermediate disease may be asymptomatic, while advanced disease can substantially impair reading, driving, facial recognition, and other daily activities.  

\subsection*{Pathophysiology}

AMD is considered a multifactorial disease driven by interactions among aging, genetic susceptibility, environmental exposures, inflammation, and metabolic dysfunction.

Major biological processes implicated include chronic inflammation, complement pathway dysregulation, oxidative stress, lipid and lipoprotein accumulation, impaired extracellular matrix maintenance and choriocapillaris degeneration.
Normal aging causes reduced photoreceptor density lipid deposition within Bruch membrane reduced density of choriocapillaris vessels and increased vascular resistance

These alterations disrupt retinal homeostasis and contribute to extracellular deposit formation. The hallmark lesion of AMD is the druse (plural: drusen), consisting of lipids, proteins, and minerals accumulating between the RPE and Bruch membrane.  

\subsection*{Classification of AMD}

\begin{table}[H]
\centering
\footnotesize
\renewcommand{\arraystretch}{1.35}

\begin{tabularx}{\textwidth}{>{\raggedright\arraybackslash}p{2.4cm}
>{\raggedright\arraybackslash}p{3.0cm}
>{\raggedright\arraybackslash}p{3.0cm}
>{\raggedright\arraybackslash}p{3.2cm}
>{\raggedright\arraybackslash}X}
\toprule
\textbf{Stage} & \textbf{Defined By} & \textbf{Clinical Characteristics} & \textbf{OCT Findings} & \textbf{Notes} \\
\midrule

\textbf{Early AMD}
&
Medium drusen ($>63\,\mu m$ and $\leq 125\,\mu m$)\newline
No AMD-associated pigmentary abnormalities
&
Usually asymptomatic\newline
Visual acuity often preserved
&
Small dome-shaped RPE elevations corresponding to drusen\newline
Relatively preserved retinal architecture
&
Progression risk is relatively low at this stage.
\\
\midrule

\textbf{Intermediate AMD}
&
Large drusen ($>125\,\mu m$)\newline
And/or pigmentary abnormalities
&
Difficulty in low-light conditions\newline
Reduced contrast sensitivity\newline
Metamorphopsia\newline
Mild visual blurring
&
Larger drusen elevations\newline
Hyperreflective foci\newline
Pigment migration\newline
Structural irregularities of the RPE
&
Most important stage for prediction studies because progression to advanced disease frequently occurs from this state.
\\
\midrule

\textbf{Late AMD: Geographic Atrophy}
&
Confluent regions of photoreceptor loss\newline
RPE loss\newline
Choriocapillaris degeneration
&
Progressive central vision loss\newline
Central scotomas\newline
Declining visual acuity
&
Loss of outer retinal layers\newline
RPE attenuation or absence\newline
Increased signal transmission into the choroid\newline
Hypertransmission defects
&
Geographic atrophy enlarges over time and is associated with irreversible visual loss.
\\
\midrule

\textbf{Late AMD: Neovascular}
&
Growth of abnormal blood vessels beneath or within the retina\newline
VEGF-mediated angiogenesis
&
Often rapid symptom onset\newline
Sudden vision decline\newline
Metamorphopsia\newline
Central visual defects
&
Intraretinal fluid\newline
Subretinal fluid\newline
Pigment epithelial detachments\newline
Hemorrhage-associated abnormalities
&
Responsible for most severe acute vision loss associated with AMD.
\\

\bottomrule
\end{tabularx}

\caption{Clinical classification of age-related macular degeneration (AMD).}
\label{tab:amd_classification}
\end{table}

\subsection*{Major Risk Factors}

\subsubsection*{Age}

Age is the strongest risk factor.

The incidence of late AMD increases dramatically with age:

\begin{itemize}
    \item 0.3 per 1000 individuals aged 55--59
    \item 5.7 per 1000 individuals aged 75--79
    \item 36.7 per 1000 individuals aged 90+
\end{itemize}

\subsubsection*{Genetics}

Estimated heritability is approximately 71\%. 

Major genetic loci include CFH and ARMS2-HTRA1, which are involved in complement regulation, inflammation, lipid metabolism, and extracellular matrix maintenance.

\subsubsection*{Smoking}

Smoking is the strongest established environmental risk factor.
Studies consistently demonstrate substantially increased AMD incidence among heavy smokers.

\subsubsection*{Diet and Lifestyle}

Protective factors include fish consumption, fruits and veggies, and physical activity.

\subsection*{Features Relevant to Progression Prediction}

For a progression-prediction project, the review repeatedly highlights several structural features associated with worsening disease.

\subsubsection*{Drusen Burden}

Important characteristics include the number of drusen, total area, volume, and presence of large drusen. Increasing drusen burden characterizes progression from early to intermediate AMD.

\subsubsection*{Pigmentary Abnormalities}

Pigment changes arise from migration of RPE cells into inner retinal layers.

These appear as:

\begin{itemize}
    \item Hyperpigmentation on fundus imaging
    \item Hyperreflective foci on OCT
\end{itemize}

The review identifies pigmentary abnormalities as a major marker of progression risk.

\subsubsection*{Hyperreflective Foci}

Migration of RPE cells creates hyperreflective structures visible on OCT.

These likely reflect active retinal remodeling and degeneration.

\subsubsection*{Outer Retinal Degeneration}

Progressive loss of photoreceptors, ellipsoid zone integrity, and RPE structure may preced advanced disease development.

\subsubsection*{Hypertransmission Defects}

Loss or attenuation of the RPE permits increased OCT signal penetration into deeper tissue.

These appear as vertical transmission columns extending into the choroid.

Such findings are particularly relevant to the current project because the proposed "barcoding" phenomenon appears visually related to localized transmission abnormalities beneath regions of RPE degeneration.

\subsubsection*{Geographic Atrophy Precursors}

Potential indicators of future atrophy include RPE disruption, hypertransmission, outer retinal thinning, or progressive drusen remodeling.

\subsubsection*{High-Risk Clinical Profile}

The AREDS study reported that:

\begin{itemize}
    \item 25.9\% of individuals with large drusen and pigment abnormalities progressed to late AMD within 5 years
    \item If one eye already had advanced AMD, progression in the fellow eye reached 53.1\% within 5 years when large drusen and pigment abnormalities were present
\end{itemize}

These findings make large drusen and pigmentary abnormalities particularly important baseline predictors.  

\section{The Barcode Sign in Optical Coherence Tomography}

\subsection*{Overview}

El Zawahry et al. described a characteristic OCT imaging pattern termed the \emph{barcode sign}, observed in eyes with extensive corneal neovascularization.
The barcode sign appears as multiple vertically oriented dark stripes extending through the OCT image. These stripes arise from attenuation of the OCT signal caused by structures that strongly block light transmission.
Although the study focuses on corneal vascularization rather than age-related macular degeneration (AMD), it provides an important example of how biological structures can generate stripe-like OCT patterns through optical shadowing mechanisms.

\subsection*{Study Design}

The authors examined five patients with extensive corneal neovascularization arising from:

\begin{itemize}
    \item Infectious keratitis
    \item Limbal stem cell deficiency
    \item Interstitial keratitis
\end{itemize}

Anterior-segment OCT scans were obtained using a Heidelberg Spectralis OCT system.

For each patient:

\begin{itemize}
    \item Twenty-one scans were acquired across the superior cornea
    \item Twenty-one scans were acquired across the inferior cornea
    \item Forty-two scans were analyzed per eye
\end{itemize}

OCT findings were compared against slit-lamp examination and clinical imaging.

\subsection*{Description of the Barcode Sign}

Major vascular trunks within the cornea completely obstructed OCT light transmission.

This produced:

\begin{itemize}
    \item A vertical hypo-reflective stripe
    \item Extension of the stripe posterior to the vessel
    \item Stripe width proportional to vessel width
    \item Multiple parallel stripes when several vessels were present
\end{itemize}

When many vessels were visualized simultaneously, the OCT image resembled a retail barcode pattern, leading to the term "Barcode Sign" where the stripes appeared as dark linear regions embedded within the brighter OCT background.

\subsection*{Proposed Mechanism}

The authors hypothesized that the shadowing effect arises primarily from the blood column within actively perfused vessels rather than the vessel wall itself.

Evidence supporting this interpretation includes:

\begin{itemize}
    \item Both afferent and efferent vessels generated equivalent shadows
    \item Shadow intensity did not appear dependent on vessel-wall structure
    \item Ghost vessels lacking active blood flow did not generate shadows
\end{itemize}

Therefore, the attenuation pattern appears linked to the presence of circulating blood.
The barcode sign may thus represent an OCT manifestation of active vascular perfusion.

\subsection*{Relationship to OCT Physics}

The barcode sign illustrates a general OCT principle: Structures that strongly attenuate or absorb OCT light can generate vertical shadow artifacts extending into deeper tissue layers.

The resulting pattern consists of:

\begin{itemize}
    \item Localized signal loss
    \item Vertical transmission defects
    \item Stripe-like image structures
    \item Reduced visibility of deeper anatomy
\end{itemize}

Such shadowing artifacts are well known in OCT imaging and can reveal information about the underlying structure causing attenuation.

\subsection*{Relevance to AMD Progression Research}

Although the barcode sign was identified in corneal vascularization, the study is conceptually relevant to AMD imaging.
Several advanced AMD features can alter OCT signal transmission, including:

\begin{itemize}
    \item Retinal pigment epithelium (RPE) degeneration
    \item Geographic atrophy
    \item Hypertransmission defects
    \item Choroidal structural abnormalities
    \item Neovascular complexes
\end{itemize}

These abnormalities can create vertically organized intensity patterns within OCT volumes.
Consequently, the barcode sign provides evidence that biologically meaningful stripe-like structures may emerge naturally from tissue-specific optical interactions rather than image-processing artifacts.

\subsection*{Implications for the Current Project}

The AMD progression project investigates whether stripe-like ``barcoding'' patterns visible within OCT scans are associated with future disease progression.

The barcode sign literature suggests several hypotheses:

\begin{itemize}
    \item Stripe patterns may arise from localized attenuation or transmission phenomena.
    \item Stripe width may reflect properties of the underlying anatomical structure.
    \item Stripe density may provide information about disease burden.
    \item Quantitative characterization of stripe morphology may produce useful imaging biomarkers.
\end{itemize}

Potential quantitative descriptors include stripe count, spacing, width, contrast, orientation, or spatial frequency content.

While the biological mechanism in AMD may differ from corneal vascular shadowing, this study provides an important precedent demonstrating that barcode-like OCT patterns can arise from meaningful underlying tissue structure and physiology.

\section{Self-Attention CNN for Retinal Layer Segmentation}

\subsection*{Overview}

Cao et al. proposed a hybrid convolutional neural network (CNN) and Transformer architecture for retinal layer segmentation in optical coherence tomography (OCT) images. The primary objective was to improve segmentation accuracy by combining the local feature extraction capabilities of CNNs with the long-range dependency modeling provided by self-attention mechanisms.

The proposed architecture extends the U-Net framework by incorporating:

\begin{itemize}
    \item A Transformer-based self-attention module,
    \item Attention Gates (AGs),
    \item Channel attention mechanisms,
    \item A one-dimensional attention strategy tailored to retinal OCT anatomy.
\end{itemize}

The model was evaluated on public retinal OCT datasets and demonstrated improved segmentation performance compared to several existing CNN-based and Transformer-based approaches.

\subsection*{Motivation}

Retinal OCT images possess highly structured anatomical organization.
Traditional CNNs perform well at extracting local image features but have limited receptive fields, making it difficult to capture long-range spatial relationships.

Challenges in OCT segmentation include layer boundaris spanning large portions of the image, anatomical variability, pathology-induced distortion, and loss of contextual information during downsampling.

The authors argue that successful segmentation requires both local texture information and global structural context.
Self-attention mechanisms are introduced to address this limitation.

\subsection*{Self-Attention Mechanism}

The Transformer component allows image regions to interact with one another regardless of spatial distance.
Self-attention is computed as
\[\text{Attention}(Q,K,V)=\text{softmax}\left(\frac{QK^\top}{\sqrt d}\right)V,\]

where $Q$ = query matrix, $K$ = key matrix, $V$ = value matrix, and $d$ = feature dimension.

Unlike convolution, which uses local neighborhoods, self-attention provides a global receptive field.
This allows the model to learn relationships between distant retinal structures and improve segmentation consistency across the image.

\subsection*{One-Dimensional Transformer Design}

A key innovation of the paper is the use of a one-dimensional Transformer.
The authors observed that retinal anatomy is organized primarily along the vertical direction.
Instead of computing attention across the entire two-dimensional image, feature maps are rearranged so that attention operates only along the retinal depth dimension.

Feature maps are reshaped from $$(B,C,H,W) \rightarrow (BW,H,C)$$ where $B$ = batch size, $C$ = number of channels, $H$ = image height, and $W$ = image width.

This design preserves anatomical structure, reduces computational cost, increases effective receptive field, and improves efficiency relative to full 2D attention.

\subsection*{Encoder--Decoder Architecture}

The overall architecture follows the encoder--decoder paradigm of U-Net.

\subsubsection*{Encoder}

The encoder consists of repeated blocks of:
\begin{itemize}
    \item $3\times3$ convolutions,
    \item Batch normalization,
    \item ReLU activation,
    \item Max pooling.
\end{itemize}
As depth increases, spatial resolution decreases and feature complexity increases. The encoder extracts increasingly abstract representations of retinal structures.

\subsubsection*{Decoder}

The decoder performs:

\begin{itemize}
    \item Upsampling,
    \item Feature fusion,
    \item Boundary reconstruction,
    \item Pixel-wise segmentation.
\end{itemize}

Skip connections transfer high-resolution information from encoder layers to corresponding decoder layers.
These connections help preserve fine retinal boundary details that may otherwise be lost during downsampling.

\subsection*{Attention Gates}

The model incorporates Attention Gates (AGs) within skip connections.
Rather than passing all encoder features directly to the decoder, attention mechanisms determine which features are most relevant.

Conceptually:\[\tilde{x}=\alpha x,\quad\forall \quad 0 \le \alpha \le 1\]
where $\alpha$ is a learned attention coefficient.

Benefits include regularization (suppression) of irrelevant features, improved focus on retinal boundaries, and enhanced segmentation precision.

\subsection*{Channel Attention}

In addition to spatial attention, channel attention is used to weight feature channels according to importance.
Different channels often capture different types of information:

\begin{itemize}
    \item Edges,
    \item Layer boundaries,
    \item Texture patterns,
    \item Anatomical structures.
\end{itemize}

Channel attention allows the network to emphasize informative channels and reduce the influence of less useful features.

\subsection*{Loss Function}

Training uses a weighted combination of Dice loss and multiclass cross-entropy loss.

The Dice coefficient is \[\text{Dice}=\frac{2|X\cap Y|}{|X|+|Y|}.\]

Corresponding Dice loss $L_{\text{Dice}}$ and cross-entropy loss $L_{\text{CE}}$ are:
\[L_{\text{Dice}}=1-\text{Dice}. \quad L_{\text{CE}}=-\sum_i y_i \log(p_i).\]

The overall objective is\[L=\lambda_1 L_{\text{Dice}}+\lambda_2 L_{\text{CE}}.\]

Combining the two losses balances pixel level calssification accuracy and global segmentation overlap.

\subsection*{Datasets}

The model was evaluated on two retinal OCT datasets:

\begin{itemize}
    \item DUKE Diabetic Macular Edema (DME) Dataset,
    \item Optic Disc Retinal OCT Dataset.
\end{itemize}

These datasets contain manually annotated retinal layer segmentations used as ground truth labels.

\subsection*{Results}

Performance was evaluated using the Dice coefficient. For the DUKE DME dataset, an average score of 0.871 was reported, and for the optic-disc dataset, 0.820 was reported.
This outperforms UNet, ATtention UNet, TransUNet, RelayNet, and DRUNey. The improvements suggest that combining CNN-based local feature extraction with Transformer-based global context modeling can improve retinal layer segmentation performance.

\subsection*{Key Takeaways}

The primary contribution of this work is the integration of self-attention into a U-Net style segmentation framework for retinal OCT images.

Important observations include:

\begin{itemize}
    \item CNNs effectively learn local retinal texture and boundary information.
    \item Self-attention provides global contextual information and long-range dependency modeling.
    
    \item Restricting attention to the vertical dimension leverages the layered anatomical structure of retinal OCT images.
    
    \item Attention Gates improve feature selection during decoding.
    
    \item Combining local and global representations improves segmentation accuracy relative to standard CNN architectures.
\end{itemize}

The paper demonstrates the potential value of hybrid CNN--Transformer architectures for retinal OCT analysis and segmentation tasks.

% =================================================================================
\part{Project Directions}


% =================================================================================
\part{Notes}

\section{Statistical Methods for Barcode Quantification}

\subsection*{Fourier Analysis}

The Fourier transform decomposes an image into spatial frequency components and is one of the most natural tools for analyzing repetitive stripe-like patterns.

Let $I(x,y)$ denote an OCT image. Its two-dimensional Fourier transform is

\[
F(u,v)
=
\sum_x
\sum_y
I(x,y)
e^{-2\pi i(ux+vy)}.
\]

The corresponding power spectrum is

\[
P(u,v)
=
|F(u,v)|^2.
\]

Interpretation:

\begin{itemize}
    \item Low frequencies correspond to broad intensity variations.
    \item High frequencies correspond to rapidly changing structures.
    \item Periodic stripe patterns generate concentrated spectral energy at characteristic frequencies.
\end{itemize}

For barcode-like transmission artifacts, a common approach is to collapse the choroidal region into a one-dimensional signal

\[
I(x)
=
\int_{\text{choroid}}
I(x,z)\,dz,
\]

and then analyze its frequency content.

If stripe spacing is approximately $d$ then its dominant frequency is $f=\frac{1}{d}$

A possible Barcoding Index is\[B=\sum_{f\in\mathcal F}|F(f)|^2,\]

where $\mathcal F$ is a frequency band corresponding to stripe spacing.

This measure is highly interpretable, computationally efficient, and quantifies periodicity well. The trade-off being it assumes approximate uniform spacing and therefore loses its spatial localization.

\subsection*{Gabor Filters}

Gabor filters are localized frequency detectors that simultaneously measure orientation, frequency, and spatial location.
A two-dimensional Gabor filter is\[g(x,y)=\exp\left(-\frac{x'^2+\gamma^2y'^2}{2\sigma^2}\right)\cos(2\pi f x'+\phi),\]
where\[x'=x\cos\theta+y\sin\theta,\qquad y'=-y\sin\theta+x\cos\theta.\]
and $\theta$ = orientation, $f$ = frequency, $\sigma$ = scale, $\gamma$ = aspect ratio, $\phi$ = phase shift.
The response of the filter is\[R(x,y)=(I*g)(x,y).\]
Strong responses occur when image structures match the filter orientation and frequency.

For OCT barcoding:

\begin{itemize}
    \item Vertical filters isolate transmission columns.
    \item Multiple frequencies allow detection of varying stripe widths.
    \item Average filter response can serve as a barcode score.
\end{itemize}

Unlike Fourier analysis, Gabor filters preserve spatial localization.

\subsection*{Structure Tensor}

The structure tensor is a classical method for quantifying local orientation.

First compute image gradients\[I_x=\frac{\partial I}{\partial x},\qquad I_y=\frac{\partial I}{\partial y}.\]
The local structure tensor is\[J=
\begin{bmatrix}
I_x^2 & I_xI_y \\
I_xI_y & I_y^2
\end{bmatrix}.
\]
In practice, local averaging is performed:\[S=G_\sigma * J,\]
where $G_\sigma$ is a Gaussian smoothing kernel.

Let $\lambda_i$ be the eigenvalues of $S$.

Interpretation:

\begin{itemize}
    \item $\lambda_1 \approx \lambda_2$:
    isotropic texture.
    \item $\lambda_1 \gg \lambda_2$:
    strong directional structure.
\end{itemize}

The dominant eigenvector indicates the primary orientation of local image patterns.
For OCT images, structure tensors quantify the strength and orientation of vertical transmission artifacts.

\subsection*{Anisotropy Metrics}

Anisotropy measures how strongly image structure is aligned along a preferred direction.

A common metric is
\[
A
=
\frac{\lambda_1-\lambda_2}
{\lambda_1+\lambda_2}, \forall \quad A \in [0,1]
\]
Interpretation:
\begin{itemize}
    \item $A=0$:
    completely isotropic texture. (weak directional organization)
    \item $A\approx1$:
    highly directional pattern. (strong vertical stripe structure)
\end{itemize}
An anisotropy score is attractive because it is simple and scale independent, and may correlate with disease severity.

\subsection*{Latent Variable Bayesian Models}

One challenge is that ``barcoding'' is not a directly observed biological structure.
Instead, it may represent a latent imaging phenomenon associated with underlying degeneration.
Let $Z_i$ be latent barcode severity for eye $i$.

Observed image features\(X_i=(B_i,G_i,A_i,\ldots)\) and may include:
\begin{itemize}
    \item Fourier-based barcode scores
    \item Gabor filter responses
    \item Anisotropy measurements
    \item Texture features
\end{itemize}

A hierarchical model can be written as\[X_i\mid Z_i\sim p(X_i\mid Z_i),\qquad Y_i\mid Z_i\sim p(Y_i\mid Z_i),\]

where $Y_i$ is the progression outcome. Posterior inference is based on \[p(Z_i\mid X_i,Y_i)=\frac{p(Y_i\mid Z_i)p(X_i\mid Z_i)p(Z_i)}{p(X_i,Y_i)}.\]

Advantages:

\begin{itemize}
    \item Separates measurement noise from biological signal.
    \item Produces uncertainty estimates.
    \item Allows probabilistic risk prediction.
\end{itemize}

Rather than classifying scans as barcoding present or absent, the model estimates a posterior distribution over severity.

\subsection*{Using Barcode Scores as Predictors}

The ultimate objective is progression prediction. Suppose $S_i$ is a barcode score computed from image features.
A logistic regression model can be written as\[\log\left(\frac{p_i}{1-p_i}\right)=\beta_0+\beta_1 S_i+\beta_2 X_{i1}+\cdots+\beta_p X_{ip},\]
where $p_i$ is the probability of preogression.

For time-to-event outcomes, a Cox model may be used:\[h(t\mid X)=h_0(t)\exp(\beta^TX).\]

Important questions include:

\begin{itemize}
    \item Does the barcode score predict progression?
    \item Does it improve prediction beyond existing AMD biomarkers?
    \item Is the score reproducible across scanners and readers?
\end{itemize}

Ultimately, barcode quantification is useful only if it provides clinically meaningful prognostic information.

\section{CNNs for OCT Analysis}

\subsection*{Convolution Layers}

The core operation of a CNN is convolution.

Given an image $X\in\mathbb R^{H\times W}$ and a filter $K\in\mathbb R^{k\times k},$
the convolution operation is

\[
Y(i,j)
=
\sum_m
\sum_n
K(m,n)
X(i+m,j+n).
\]

Each filter learns a feature detector. Early CNN layers typically learn edges, gradients, and simple textures, whereas deeper layers learn anatomical structures, retinal layers, and disease-related patterns. Important hyperparameters are kernel size, stride, zero-padding, and number of filters, and output dimensions satisfy:
\[
W_{\text{out}}
=
\frac{W-F+2P}{S}+1.
\]

\subsection*{Pooling Layers}

Pooling reduces spatial dimensionality.

Max pooling computes

\[
Y(i,j)
=
\max_{(m,n)\in\mathcal N}
X(i+m,j+n).
\]

Average pooling computes

\[
Y(i,j)
=
\frac1{|\mathcal N|}
\sum_{(m,n)\in\mathcal N}
X(i+m,j+n).
\]

Benefits:

\begin{itemize}
    \item Reduced computation and overfitting
    \item Increased translation invariance
\end{itemize}

Typical pooling windows are either $2\times2$ or $3\times3$.

\subsection*{Activation Functions}

Activation functions introduce nonlinearity.

ReLU: $f(x)=\max(0,x).$

Leaky ReLU: $f(x)=\max(\alpha x,x).$

Sigmoid: $f(x)=\frac{1}{1+e^{-x}}.$

Tanh: $f(x)=\frac{e^x-e^{-x}}{e^x+e^{-x}}.$


ReLU is dominant in modern CNNs because it mitigates vanishing gradients and is computationally efficient.

\subsection*{Fully Connected Layers}

After convolutional feature extraction, feature maps are flattened into a vector $x\in\mathbb R^d.$ A fully connected layer computes\[y=Wx+b.\]

Classification probabilities are obtained using softmax: \[p_k=\frac{e^{z_k}}{\sum_j e^{z_j}}.\]

For binary classification:\[p=\frac{1}{1+e^{-z}}.\]

Fully connected layers combine learned image features into disease predictions.

\subsection*{CNN Feature Learning}

A key advantage of CNNs is automatic feature learning. CNNs learn feature representations directly from data.

The representation learned by layer $\ell$ can be written as $h_\ell=f_\ell(h_{\ell-1}).$
Successive layers learn increasingly abstract image representations:
\[\text{Edges}\rightarrow\text{Textures}\rightarrow\text{Shapes}\rightarrow\text{Pathology}.\]

\subsection*{ResNet}

Deep networks suffer from vanishing gradients. ResNet introduces residual connections:
\[y=F(x)+x.\]
The identity shortcut allows gradients to flow through many layers. This is also easier to optimize and has better convergence, and most modern OCT classifiers use the ResNet backbone (EX: ResNet-18/34/50/101).

\subsection*{U-Net}

U-Net is the dominant architecture for biomedical image segmentation.

Architecture:

\begin{itemize}
    \item Encoder (downsampling)
    \item Bottleneck
    \item Decoder (upsampling)
    \item Skip connections
\end{itemize}

The encoder extracts context while the decoder restores spatial resolution. Skip connections preserve fine anatomical detail.

\subsection*{Transfer Learning}

Medical imaging datasets are often small. Transfer learning uses a model pretrained on a large dataset (e.g., ImageNet) and adapts it to a new task.

Let $f_\theta$ denote a pretrained CNN. Training proceeds by:
\begin{enumerate}
    \item Initializing weights from the pretrained model.
    \item Replacing the final classification layer.
    \item Fine-tuning on OCT images.
\end{enumerate}

This has lower training time and generalizes better, and performs better on small datasets.

\subsection*{Interpretable CNNs (Grad-CAM)}

A common criticism of CNNs is their lack of interpretability. Gradient-weighted Class Activation Mapping (Grad-CAM) identifies image regions responsible for predictions.
Let $A^k$ denote feature map $k$. Importance weights are computed as\[\alpha_k=\frac{1}{Z}\sum_i\sum_j\frac{\partial y}{\partial A^k_{ij}},\]
where $y$ is the model output. The Grad-CAM heatmap is\[L=\text{ReLU}\left(\sum_k\alpha_kA^k\right).\]

The resulting map highlights image regions most influential for the prediction. In OCT analysis, Grad-CAM can help determine whether the network is utilizing:

\begin{itemize}
    \item RPE abnormalities
    \item Hypertransmission regions
    \item Barcode-like artifacts
    \item Other retinal structures
\end{itemize}

thereby improving interpretability and clinical trust.


\end{document}
